\section{Results}

The completed results below are deterministic/offline quality and grounding
evidence for the current PaperToSkill artifacts, plus first-pass real-reuse
execution evidence. AIDE-T1 and AIDE-T2 are scored for one GPT-family run on a
locked Kaggle Spaceship Titanic validation split using the same saved model
outputs and an extended 300-second local scorer budget. AIDE-T1 is solved by
both Summary and PaperToSkill with nearly tied validation scores (0.816 and
0.817), while AIDE-T2 is a PaperToSkill-only success (0.826 versus a Summary
timeout). SWE-T1 is scored for one GPT-family run on a locked
SWE-Bench Lite SQLFluff instance: both Summary and PaperToSkill produce patches
that fail to apply and score 0.000. SWE-T2 is scored for one GPT-family run on
a locked SWE-Bench Verified-style Astropy instance: the Summary patch fails to
apply and scores 0.000, while the PaperToSkill patch applies, passes the hidden
target tests, and scores 1.000. In the real-reuse table, REF-T1 and REF-T2 are
also scored for one GPT-family run: both Summary and PaperToSkill reach 1.000
on the locked Reflexion QA and retry fixtures. SNAP-T1 and SNAP-T2 are scored
for one GPT-family run over prepared official miniature fixtures. PaperToSkill
scores above Summary on these two SnapATAC2 dry-run tasks (0.500 vs. 0.000 and
0.400 vs. 0.200), but all SNAP scores remain below the success threshold,
mainly because the responses did not produce complete runtime, memory, and
quality artifacts. Across the eight filled tasks, the runner, scorer, and
table-update paths are live. The outcome is mixed: AIDE-T2 and SWE-T2 are
positive PaperToSkill rows, AIDE-T1 and REF show little or no advantage because
both conditions solve the fixtures, and SWE-T1/SNAP primarily expose
patch-application and artifact-completion failure boundaries. These results do
not establish
aggregate advantage over Summary. Table~\ref{tab:real-reuse-failure-analysis}
summarizes the row-level boundary modes and the method contracts they suggest
for follow-up work. The SWE-T1 diagnostics separate two contract issues. The
shared-source-context follow-up
(Table~\ref{tab:swe-t1-source-context-followup}) fixes patch application for
both conditions but still fails the hidden target test. The issue-aligned
follow-up (Table~\ref{tab:swe-t1-issue-aligned-followup}) then uses a
pre-registered revised scorer and both Summary and PaperToSkill pass, so it
shows that the revised task contract can close rather than showing a
PaperToSkill advantage. A paired SnapATAC2 executable-artifact diagnostic
(Table~\ref{tab:snapatac2-executable-followup}) then runs a pre-registered
controlled scaffold over the same miniature fixtures. It reaches the SNAP
scorer's artifact/runtime/memory contract for both Summary and PaperToSkill
conditions, confirming that the prior SNAP failure mode is an execution
contract boundary rather than a fixture or scorer impossibility; it is not a
main-row replacement and does not show a PaperToSkill advantage. A subsequent
model-generated executable-candidate diagnostic for SNAP-T1
(Table~\ref{tab:snapatac2-executable-candidate-followup}) also reaches 1.000
for both Summary and PaperToSkill under the pre-registered executable-candidate
contract. This closes the candidate-script artifact contract for both
conditions; it does not replace the main SNAP rows or show a PaperToSkill
advantage. The auxiliary Full Excerpt sanity subset does not reverse this
interpretation: AIDE-T1 and SWE-T1 score 0.000 under Full Excerpt, while
SNAP-T1 scores 0.250, below both PaperToSkill and the success threshold.
All four generated skills score 20/20 on deterministic skill rubrics. On context
coverage, the generated skills score 7.867/9 for AI Scientist-v2, 8.267/9 for
Reflexion, 9.1/10 for AIDE, and 8.9/10 for Toolformer. The generic-summary and
abstract-only baselines score substantially lower across all four cases
(Table~\ref{tab:main-results}).

The generated skills remain compact: 782 words for AI Scientist-v2, 479 words
for Reflexion, 927 words for AIDE, and 943 words for Toolformer, all under the
1200-word budget. Source validation finds support rates of 0.938, 1.0, 1.0,
and 1.0 with zero invalid line ranges. The one weak AI Scientist-v2 span is
recorded as a boundary case rather than treated as fully verified.

The tokenizer-aware context proxy shows the same compactness pattern relative
to full paper contexts. Under \texttt{o200k\_base}, generated skills use between
2.39\% and 9.65\% of the input tokens required by full extracted paper text,
with the full tokenizer table retained in the package artifacts. A separate
character-proxy report preserves the earlier
$\lceil\mathrm{characters}/4\rceil$ estimate for reproducibility. We interpret
both as coverage-preserving compression relative to full-paper context, not as
provider billing evidence or a claim that skills are the shortest possible
context.

Transfer-note ablations show a consistent offline readiness pattern. Full
skills score 10/10 on the readiness metric across all four curated papers.
Removing \texttt{Transfer Notes} lowers readiness to 7.6/10 in all four cases,
while generic summaries score between 1.2/10 and 2.25/10. This supports the
narrower claim that transfer notes encode artifact-readiness signals. It does
not yet prove improved live cross-harness success.

The live-transfer saved-response evaluation covers all four paper packets.
AI Scientist-v2, Reflexion, AIDE, and Toolformer each have six saved responses
across Codex-style and Claude-style harness prompts and three context variants.
The aggregate evaluator reports 24 total rows, 24 scored rows, 0 pending rows,
and average normalized score 1.0. AI Scientist-v2, Reflexion, and AIDE rows
score 11/11 each; Toolformer rows score 9/9 each. The AIDE run includes one
provider fallback row where \texttt{claude-opus-4-8} closed the connection and
\texttt{claude-opus-4-7} succeeded. These are deterministic output-contract
scores over saved response files, not human-rated semantic fidelity or
downstream execution-outcome rates.

The automatic note scaffold is evaluated separately on Toolformer and AIDE. The
Toolformer auto-note-derived skill scores 20/20 on the deterministic rubric,
9.3/10 on context coverage, 10/10 on offline transfer readiness, and 1.0
source-span support with zero invalid ranges. The AIDE auto-note-derived skill
scores 20/20, 8.467/10 on context coverage, 9.5/10 on offline transfer
readiness, and 1.0 source-span support with zero invalid ranges. The full
auto-note comparison table is retained in the package artifacts. These results
support the narrower claim that extracted text can seed auditable note
scaffolds. The local pipeline also
includes a direct-PDF smoke path through \texttt{pdftotext -layout} that records
the generated text source in the manifest, but this does not establish robust
arbitrary-PDF automation.

We additionally compare PaperToSkill against Paper2Agent at the artifact and
workflow level using the Paper2Agent paper text and current PaperToSkill
artifacts. The bounded comparison covers required inputs, generated artifact
type, setup burden, validation checks, failure handling, source traceability,
and runtime dependency. All seven criteria are source-backed and ready in the
local comparison table. The result clarifies positioning: Paper2Agent produces
MCP servers for papers with associated codebases, while PaperToSkill produces
portable natural-language skills with explicit source and failure boundaries.
This comparison does not run Paper2Agent, deploy an MCP server, or establish an
end-to-end baseline performance result.

\section{Usage Examples and Model Ablations}

The repository includes usage examples for three practical workflows:
loading a generated skill in a Codex-style harness, generating an automatic
note scaffold from extracted text or a smoke-tested PDF source, and building
model-ablation prompts. The
model-ablation protocol produces a prompt grid for Claude Opus 4.8, a
GPT-family slot requested as GPT 5.5, and a DeepSeek slot. In the current
protocol, the saved-response evaluator reports 6 total rows, 6 scored rows, 0
pending rows, and average normalized score 1.0. The GPT-family protocol refresh
completed both rows with \texttt{gpt-5.5}. The DeepSeek run completed both rows
with \texttt{deepseek-v4-flash}. The latest Claude protocol refresh used the
correct Anthropic Messages path but was blocked by provider HTTP 502; the
scored Claude rows come from previously saved response files. A local
output-token proxy over all six saved model-ablation responses counts 9,594
\texttt{o200k\_base} output tokens. This 6/6 saved-response score is not live
downstream task success, provider billing, live-invoice, success-per-dollar
evidence, or a broad model-quality result.

Separately, the real-reuse LLM ablation attaches model variation to the
paper-task protocol rather than to the older saved-response usage-plan
protocol (Table~\ref{tab:real-reuse-llm-ablation}). The pre-registered slice
uses AIDE-T2, SWE-T2, and REF-T2 under the same Summary/PaperToSkill pairing.
GPT-family and DeepSeek-family rows are collected, while Claude-family rows
remain pending after provider HTTP 502 responses. The collected rows are mixed:
GPT-family favors Summary on AIDE-T2 and ties on REF-T2, while both conditions
fail SWE-T2; DeepSeek ties on the collected slices. This table is auxiliary
model-slice evidence and does not replace the main real-reuse rows.

The bounded AI-Scientist-v2 integration path is complete for the local marker
smoke and one full live run. The smoke report records a complete marker-contract
response, and the full-run handoff records one completion directory. The
AI-Scientist-v2-generated synthetic benchmark ties skill and full-excerpt
task-success rates at 0.80 while the skill uses fewer tokens (86.2 versus
113.2); abstract, generic-summary, and no-context baselines score 0.20, 0.00,
and 0.00. A retrieval-depth sensitivity branch reaches 1.00 only when K=all.
This is bounded integration and synthetic sensitivity evidence, not human
fidelity, real-data validation, or broad live research-task success.

\section{Discussion}

The main result is not that PaperToSkill ``understands'' papers in the broadest
sense. The result is that a structured paper-note-to-skill pipeline can preserve
operational knowledge that summaries discard. The skills include workflow
steps, validation checks, failure cases, and source anchors, which are exactly
the parts needed when a user wants an agent to apply a paper rather than merely
explain it.

The benchmark also reveals useful failure modes. During the AIDE case, the
initial extractor capped workflow, validation, and failure candidates too
aggressively, which dropped data-preview and LLM-cost content. Earlier phases
also exposed title parsing, source-audit mapping, and source-span line-counting
bugs. These records are the type of failure branches PaperToSkill should
preserve: not just final successes, but the ways extraction, evaluation, or
external dependencies can lose important operational content.

The real-reuse rows add a harsher version of the same lesson. The AIDE rows
show that task slice and local scorer budget can change the apparent boundary:
with a longer scorer budget AIDE-T1 is solved by both conditions and AIDE-T2
becomes a PaperToSkill-only success, while the Summary AIDE-T2 timeout still
argues for explicit runtime and fallback contracts. SWE-T1 shows that
software-engineering reuse can fail before tests run if a patch cannot be
applied cleanly. The SnapATAC2 rows show that partial pipelines are not enough
when the scorer requires complete runtime, memory, and quality artifacts. These
results point to the next method pressure points: budget contracts,
patch/application constraints, required artifact manifests, and recovery
instructions should be explicit when a source paper's method is converted into
a reusable skill.

\section{Limitations}

The current work has several important limits. First, the main benchmark
converts curated notes rather than arbitrary PDFs; automatic note scaffolds have
only been retained on Toolformer and AIDE extracted text. Second, the metrics
are deterministic and partly lexical, so they can over-credit exact matches or
under-credit valid paraphrases. Third, live-transfer saved-response coverage is
complete for the current prompt-packet protocol, but the scorer is
deterministic and contract-based rather than a human semantic-fidelity or live
execution-outcome evaluation. Fourth, the bounded AI-Scientist-v2 smoke and
full live run is complete, but its positive result is synthetic and should not
be read as broad live task success. Fifth, the 24-cell single-reviewer internal
annotation is now scored and summarized, but it is a bounded review artifact
rather than independent or broad human validation or live task-success
evidence. Sixth,
compactness is measured by word count, deterministic character proxy, local
tokenizer-aware input-token proxy, and local output-token proxy for saved model
responses, not by provider-specific prices, live invoices, or live success per
dollar. Seventh, the real-reuse rows are single-run GPT-family measurements,
and original-paper reference scores are reported references unless the same
data, input/output contract, metric, budget, and runtime setting are reproduced
locally. Several rows fail because generated scripts or patches do not finish
within the locked budget or do not produce complete scorer artifacts, so they
should be read as boundary evidence rather than a broad downstream
effectiveness claim.
